% =====================================================================
% page-layout-fancyhdr.tex — Running headers/footers wired via fancyhdr
%
% Maps the package-agnostic \keystone* content macros (defined by the
% page-furniture.lua filter, which owns slot composition) onto fancyhdr's
% parity-aware slot positions. The KOMA-Script sibling
% (page-layout-scrlayer.tex) maps the same macros onto
% scrlayer-scrpage's native slot API.
%
% It also owns the head-rule default for this package: the rule is off
% by default, and the \keystoneheaderrule shim (defined below, emitted
% by page-furniture.lua only when the author sets header-rule: enabled)
% re-enables it through fancyhdr's \headrulewidth.
%
% fancyhdr slot naming: the third character is the column
% (L/C/R = left/center/right), the fourth is the page parity
% (O/E = odd/even = recto/verso). Under \documentclass{article}
% or \documentclass{report} (single-sided by default), only the
% O-suffixed slots are selected; the E-suffixed wirings compile
% but never render. Two-sided documents (book / scrbook, or any
% target with classoption: twoside) get parity-distinct output.
%
% Opening pages — every page the class routes through `plain`: each
% chapter's first page, the table of contents, the index, and part
% titles — carry no running header and no rule; the heading is the
% page's only landmark. LaTeX's own `plain` already prints a centered
% folio there and nothing else, so redefining it below is what puts that
% line under the author's control: it becomes the \keystoneopeningfooter
% slots, holding `footer-text-opening`. They carry the same parity grid as
% the body rows, because an author can move that line to the outside edge
% of a spread, where recto and verso part ways.
%
% The title page is the one `plain` transition kept bare; the \maketitle
% hook below holds that for classes setting the title in-page.
%
% Auto-inserted filler blanks — the verso `\cleardoublepage` inserts
% before a recto chapter opening — are likewise kept bare, via the
% `\cleardoublepage` redefinition near the foot of this file. fancyhdr
% would otherwise print a stale running mark and the folio on them; the
% scrlayer/KOMA targets style such blanks empty by default, so this
% restores parity.
%
% This file is part of the Keystone project layout.
% =====================================================================
\usepackage{fancyhdr}
\pagestyle{fancy}

% fancyhdr's \pagestyle{fancy} would otherwise draw its 0.4pt default
% headrule on every body page; zeroing it here matches the scrlayer
% default (headsepline=false), re-raised to 0.4pt by \keystoneheaderrule
% (below). The plain pagestyle (opening pages) keeps 0pt regardless.
\renewcommand{\headrulewidth}{0pt}
\newcommand{\keystoneheaderrule}{\renewcommand{\headrulewidth}{0.4pt}}

% Dynamic-placeholder shims. The page-furniture filter emits these
% macro references when an author writes `{page}`, `{chapter}`, or
% `{section}` in `header-text` / `footer-text` content. Resolution
% is per-package: fancyhdr's marks (\leftmark, \rightmark) are
% uppercased by default, so `\nouppercase{...}` preserves the case
% the author set at the `\markboth` boundary. `\providecommand` (not
% `\newcommand`) lets a future author or test override the shim
% per-document if needed.
\providecommand{\keystonepagemark}{\thepage}
\providecommand{\keystonechaptermark}{\nouppercase{\leftmark}}
\providecommand{\keystonesectionmark}{\nouppercase{\rightmark}}
\fancyhead[LO]{\keystoneheaderLO}\fancyhead[LE]{\keystoneheaderLE}
\fancyhead[CO]{\keystoneheaderCO}\fancyhead[CE]{\keystoneheaderCE}
\fancyhead[RO]{\keystoneheaderRO}\fancyhead[RE]{\keystoneheaderRE}
\fancyfoot[LO]{\keystonefooterLO}\fancyfoot[LE]{\keystonefooterLE}
\fancyfoot[CO]{\keystonefooterCO}\fancyfoot[CE]{\keystonefooterCE}
\fancyfoot[RO]{\keystonefooterRO}\fancyfoot[RE]{\keystonefooterRE}
\fancypagestyle{plain}{%
  \fancyhf{}%
  \fancyfoot[LO]{\keystoneopeningfooterLO}\fancyfoot[LE]{\keystoneopeningfooterLE}%
  \fancyfoot[CO]{\keystoneopeningfooterCO}\fancyfoot[CE]{\keystoneopeningfooterCE}%
  \fancyfoot[RO]{\keystoneopeningfooterRO}\fancyfoot[RE]{\keystoneopeningfooterRE}%
  \renewcommand{\headrulewidth}{0pt}%
}

% Hold the "title page is not an opening page" invariant (page-furniture.lua's
% macro contract) now that `plain` carries a footer. book and report set the
% title inside a `titlepage` environment, already `empty`. article sets it
% in-page, leaving `plain` on that page — which would print the opening footer
% against a title. `\if@titlepage` is the class's own switch between those two
% shapes, so the hook lands only on the in-page one.
\makeatletter
\if@titlepage\else
  \AddToHook{cmd/maketitle/after}{\thispagestyle{empty}}
\fi
\makeatother

% Filler blanks inserted by `\cleardoublepage` — the verso it pushes before
% a recto-opening page in a two-sided document (book by default; report or
% article only with `classoption: twoside`) — inherit the active
% `\pagestyle{fancy}`, so fancyhdr prints the running head, foot, and folio
% on a page meant to be blank (e.g. a stale "Contents" mark after the TOC).
% Redefine `\cleardoublepage` to stamp `\thispagestyle{empty}` on the blank,
% matching the scrlayer/KOMA targets and the emptypage package's idiom.
%
% `\if@twoside` / `\if@twocolumn` are the class's own guards: one-sided
% documents insert no filler blank (their `\cleardoublepage` is a plain
% `\clearpage`), so preserving them leaves that case untouched. The page
% counter is not altered — only the printed folio on the blank disappears.
%
% Suppressing the folio too (not just head/foot) is deliberate: a filler
% blank is a page the reader is meant to skip over, so a folio there is
% the one piece of furniture that draws the eye back to it. KOMA
% suppresses it likewise, which is what keeps the two targets identical
% here. Note this is the one page class the opening footer does not
% reach — an author who wants folios on blanks needs a bespoke pagestyle
% (empty head, folio-only foot), not a change here.
%
% Body mirrors the class definition of `\cleardoublepage` (classes.dtx)
% verbatim, plus the `\thispagestyle{empty}` line. On a TeX Live bump,
% diff against that definition if blank-page behavior shifts.
\makeatletter
\renewcommand{\cleardoublepage}{%
  \clearpage%
  \if@twoside%
    \ifodd\c@page\else%
      \thispagestyle{empty}%
      \hbox{}\newpage%
      \if@twocolumn\hbox{}\newpage\fi%
    \fi%
  \fi%
}
\makeatother
